\documentclass[12pt]{article}

\usepackage[margin=1in]{geometry}
\usepackage{amsmath,amsthm,amssymb}
\usepackage{fancyhdr}
\pagestyle{fancy}
\usepackage{subfig}
\usepackage{float}
\usepackage[T1]{fontenc}
\usepackage[utf8]{inputenc}
\usepackage[polish]{babel}
\usepackage{graphicx}
\usepackage{tikz}
\usepackage{pgfplots}
\usepackage{amsmath}
\usepackage[mathscr]{euscript}
\let\euscr\mathscr \let\mathscr\relax
\usepackage[scr]{rsfso}
\usepackage{amsthm}
\usepackage{amssymb}
\usepackage{multicol}
\usepackage[colorlinks=true, pdfstartview=FitV, linkcolor=blue,
citecolor=blue, urlcolor=blue]{hyperref}
\usepackage{minted}
\setminted{fontsize=\footnotesize, linenos, breaklines, frame=single}

\DeclareMathOperator{\arcsec}{arcsec}
\DeclareMathOperator{\arccot}{arccot}
\DeclareMathOperator{\arccsc}{arccsc}
\newcommand{\ddx}{\frac{d}{dx}}
\newcommand{\dfdx}{\frac{df}{dx}}
\newcommand{\ddxp}[1]{\frac{d}{dx}\left( #1 \right)}
\newcommand{\dydx}{\frac{dy}{dx}}
\let\ds\displaystyle
\newcommand{\intx}[1]{\int #1 \, dx}
\newcommand{\intt}[1]{\int #1 \, dt}
\newcommand{\defint}[3]{\int_{#1}^{#2} #3 \, dx}
\newcommand{\imp}{\Rightarrow}
\newcommand{\un}{\cup}
\newcommand{\inter}{\cap}
\newcommand{\ps}{\mathscr{P}}
\newcommand{\set}[1]{\left\{ #1 \right\}}
\newtheorem*{sol}{Solution}
\newtheorem*{claim}{Claim}
\newtheorem{problem}{Problem}
\begin{document}

\lhead{Laboratorium 5}
\chead{Igor Józefowicz 193257 \\ Kamil Śliwiński 193740}
\rhead{\today}
\title{Uczenie głębokie -- Laboratorium 1 }
\author{Igor Józefowicz 193257\and Kamil Śliwiński 193740}
\maketitle
%==================================================================================


\section{zadanie 1}

\begin{minted}[frame=lines, linenos, breaklines]{python}

## Wynik dla jądra k1:
[[-4.,  -6.,   4.],

 [-20., -14.,  20.],

 [-20., -10.,  20.]]
 ## Wynik dla jądra k2:
[[ 9.,  6., 16.],

 [21., 27., 30.],

 [14., 30., 27.]]
 ## Wynik dla jądra k3:
 [[  8.,  15.,   9.],

 [  7.,   8.,   9.],

 [ -8., -15.,  -9.]]

\end{minted}



\section{zadanie 2}


\begin{minted}[frame=lines, linenos, breaklines]{python}

import matplotlib.pyplot as plt
import os

kernels['k4'] =   np.array([
                            [1., 0., 0.,0., 0.],
                            [0., 1., 0.,0., 0.],
                            [0., 0., 1.,0., 0.],
                            [0., 0., 0.,1., 0.],
                            [0., 0., 0.,0., 1.]
                            ])



image_path = 'lab1_1.jpg'
if os.path.exists(image_path):
    img_input = read_image(image_path, 200, 200)
else:
    print(f"Plik {image_path} nie został znaleziony! Używam losowego obrazu testowego.")
    img_input = np.random.rand(200, 200, 1) * 255

# Obliczenie wyników splotów
out_k1 = apply_conv2d(img_input, kernels['k1'])
out_k2 = apply_conv2d(img_input, kernels['k2'])
out_k3 = apply_conv2d(img_input, kernels['k3'])
out_k4 = apply_conv2d(img_input, kernels['k4'])

# Wyświetlenie obrazów
fig, axs = plt.subplots(1, 5, figsize=(15, 5))

# Oryginalny obraz
axs[0].imshow(img_input.squeeze(), cmap='gray')
axs[0].set_title('Oryginalny obraz')
axs[0].axis('off')

# Wynik dla k1
axs[1].imshow(out_k1[0, :, :, 0], cmap='gray')
axs[1].set_title('Filtr k1 (pionowy)')
axs[1].axis('off')

# Wynik dla k2
axs[2].imshow(out_k2[0, :, :, 0], cmap='gray')
axs[2].set_title('Filtr k2 (wyostrzenie)')
axs[2].axis('off')

# Wynik dla k3
axs[3].imshow(out_k3[0, :, :, 0], cmap='gray')
axs[3].set_title('Filtr k3 (poziomy)')
axs[3].axis('off')

# Wynik dla k4
axs[4].imshow(out_k4[0, :, :, 0], cmap='gray')
axs[4].set_title('Filtr k4')
axs[4].axis('off')

plt.tight_layout()
plt.show()

\end{minted}

\begin{figure}[H]
    \centering
    \includegraphics[width=0.65\textwidth]{p1.PNG}
    \caption{}
\end{figure}

\begin{figure}[H]
    \centering
    \includegraphics[width=0.9\textwidth]{p2.PNG}
    \caption{}
\end{figure}
\begin{figure}[H]
    \centering
    \includegraphics[width=0.9\textwidth]{p3.PNG}
    \caption{}
\end{figure}


\begin{figure}[H]
    \centering
    \includegraphics[width=0.6\textwidth]{p4.PNG}
    \caption{Kształt wejściowy obrazu: (200, 200, 1)
Kształt wyjściowy po splotu ze strides=(2,2) i padding='VALID': (1, 99, 99, 1)}
\end{figure}
 Stride powoduje zmiane rozdzielczości przetarzanego obrazu natomiast padding ustawiony na 'VALID 'nie stosuje żadnego dopełnienia (obraz wyjściowy będzie mniejszy, ponieważ filtr nakłada się tylko tam, gdzie w pełni mieści się na obrazie)
\section{zadanie 3}

\textbf{Dane wejściowe:} \\
Wymiar początkowy (krok 0): $24 \times 24 \times 1$ dla \texttt{input = np.ones((24, 24, 1))}

\vspace{0.5cm}
\textbf{Przebieg operacji:}

\begin{enumerate}
    \item \textbf{Splot (Conv2D):} 16 filtrów o rozmiarze $3 \times 3$, krok $(1,1)$, padding \texttt{SAME}
    \begin{itemize}
        \item \textbf{Kształt wyjściowy:} $24 \times 24 \times 16$
    \end{itemize}

    \item \textbf{Pooling:} jądro $5 \times 5$, krok $(3,3)$, padding \texttt{VALID}
    \begin{itemize}
        \item \textbf{Kształt wyjściowy:} $7 \times 7 \times 16$
        \item \textit{Uwaga:} Wymiar przestrzenny wynosi $7 \times 7$, a nie $8 \times 8$, ponieważ padding \texttt{VALID} oznacza brak marginesów, co prowadzi do ucięcia pikseli niemieszczących się w pełnym oknie. Wynika to ze wzoru: 
        $$W_{wyj} = \left\lfloor \frac{24 - 5}{3} \right\rfloor + 1 = 6 + 1 = 7$$
    \end{itemize}

    \item \textbf{Splot (Conv2D):} 32 filtry o rozmiarze $3 \times 3$, krok $(1,1)$, padding \texttt{SAME}
    \begin{itemize}
        \item \textbf{Kształt wyjściowy:} $7 \times 7 \times 32$
    \end{itemize}

    \item \textbf{Pooling:} jądro $3 \times 3$, krok $(2,2)$, padding \texttt{VALID}
    \begin{itemize}
        \item \textbf{Kształt wyjściowy:} $3 \times 3 \times 32$
        \item \textit{Obliczenie:} 
        $$W_{wyj} = \left\lfloor \frac{7 - 3}{2} \right\rfloor + 1 = 2 + 1 = 3$$
    \end{itemize}
\end{enumerate}

\vspace{0.5cm}
\textbf{Ostateczny wymiar macierzy:}
$$3 \times 3 \times 32$$



\section{zadanie 4}

\subsection{zadanie 4a}

\begin{figure}[H]
    \centering
    \includegraphics[width=0.9\textwidth]{4a_model_summary.png}
    \caption{}
\end{figure}

\begin{figure}[H]
    \centering
    \includegraphics[width=0.9\textwidth]{4a_training_plot.png}
    \caption{Wykres przedstawia proces treningu.}
\end{figure}

Wynik ewaluacji - dokładność (accuracy):  0.48170000314712524

\subsection{zadanie 4b}

\begin{figure}[H]
    \centering
    \includegraphics[width=0.9\textwidth]{4b_model_summary.png}
    \caption{}
\end{figure}

\begin{figure}[H]
    \centering
    \includegraphics[width=0.9\textwidth]{4b_training_plot.png}
    \caption{Wykres przedstawia proces treningu.}
\end{figure}
Wynik ewaluacji - dokładność (accuracy):  0.6442000269889832

\begin{figure}[H]
    \centering
    \includegraphics[width=0.9\textwidth]{4b_example_images.png}
    \caption{}
\end{figure}

Komentarz do wyników modelu z zadania 4a (prosty model) i zadania 4b - model rozbudowany o 1 dodatkową warstwę splotową i 2 dodatkowe max pooling 2d:

Model 4a uzyskał dokładność ~48.2\%. Model 4b poprawił ten wynik o około 16\% osiągając dokładność 64.4\%, co jest zauważalną poprawą, ale wciąż nie jest zadowalającym wynikiem.

\subsection{zadanie 4c}

\begin{figure}[H]
    \centering
    \includegraphics[width=0.6\textwidth]{4c_example_image.png}
    \caption{Przykładowe zdjęcie}
\end{figure}

\begin{figure}[H]
    \centering
    \includegraphics[width=0.9\textwidth]{4c_heatmap_filter_5.png}
    \caption{Mapa ciepła dla 3 kanałów - RGB - Red, Green, Blue}
\end{figure}

\subsection{zadanie 4d}

\begin{figure}[H]
    \centering
    \includegraphics[width=0.9\textwidth]{4d_model_summary.png}
    \caption{}
\end{figure}

\begin{figure}[H]
    \centering
    \includegraphics[width=0.9\textwidth]{4d_training_plot.png}
    \caption{Wykres przedstawia proces treningu.}
\end{figure}
Wynik ewaluacji - dokładność (accuracy):  0.7318000197410583 \\
Model z zadania 4d został rozbudowany o trzecią warstwę splotową z 256 filtrami. Po tej trzeciej warstwie konwolucyjnej dodano jeszcze warstwę max pooling 2d.

Model trenował się zauważalnie dłużej od poprzednich.

W modelu z zadania 4d zdecydowanie widać overfitting. Dokładność na zbiorze treningowym wynosi ponad 95\%, natomiast na zbiorze testowym wynosi 73.1\%. Natomiast ten wynik 73.1\% wciąż jest lepszy od modeli 4a czy 4b.  



\end{document}
